% 第六章 重点公司分析

本章对 10 家核心 A 股功率半导体标的进行深度分析，从业务结构、技术进展、AI 弹性、财务质量、估值水平、催化与风险六个维度展开。排序基于综合推荐优先级。

\section{时代电气（688187）— 估值最低的功率半导体 IDM}

\subsection{基本信息与业务概览}

时代电气是中国中车旗下上市公司，业务涵盖轨道交通装备和新兴装备两大板块。功率半导体是其新兴装备板块的核心业务，受益于新能源汽车和新能源发电快速增长，已成为公司第二增长曲线。

公司是国内少数具备 IGBT 芯片设计、制造、封装测试全流程能力的 IDM 厂商，产品覆盖 650V–6500V 全电压等级，是国内 IGBT 技术实力最强的公司之一。

\subsection{业务结构与技术进展}

\begin{exhibitbox}[表 6-1：时代电气业务结构（2025E 估算）]
\centering
\small
\begin{tabularx}{\textwidth}{L{2.5cm} X c c c}
\toprule
\textbf{业务板块} & \textbf{主要产品} & \textbf{营收占比} & \textbf{毛利率} & \textbf{增速} \\
\midrule
轨道交通装备 & 牵引系统、信号系统等 & ~55\% & ~38\% & 低个位 \\
新兴装备—功率半导体 & IGBT/SiC 器件与模块 & ~25\% & ~30-32\% & 30-40\% \\
新兴装备—工业传动 & 变频器、伺服系统等 & ~10\% & ~35\% & 5-10\% \\
新兴装备—新能源汽车电驱 & 电驱系统 & ~8\% & ~20\% & 20-30\% \\
其他 & 传感器、海工装备等 & ~2\% & — & — \\
\midrule
合计 & — & 100\% & ~35-38\% & ~15-20\% \\
\bottomrule
\end{tabularx}
\sourcenote{公司财报，本研究团队测算}
\end{exhibitbox}

技术进展亮点：
\begin{itemize}[leftmargin=2em]
\item \textbf{IGBT}：车规级 IGBT 模块已批量交付多家头部车企，是国内车规 IGBT 龙头之一
\item \textbf{SiC}：1200V SiC MOSFET 已实现车规定点并批量交付，SiC 产线持续扩产
\item \textbf{高压器件}：6500V 高压 IGBT 在轨道交通、电网领域国产替代领先
\end{itemize}

\subsection{AI/数据中心业务弹性}

时代电气在数据中心领域的布局相对较少，主要以工业级和汽车级功率器件为主。但公司的高压大功率 IGBT 和 SiC 技术储备可延伸至数据中心高压 PSU 和 PFC 场景，AI 业务弹性目前相对有限。

\subsection{财务数据与盈利能力}

\begin{exhibitbox}[表 6-2：时代电气核心财务数据]
\centering
\small
\begin{tabularx}{\textwidth}{X c c c c}
\toprule
\textbf{指标} & \textbf{2022} & \textbf{2023} & \textbf{2024} & \textbf{2025E} \\
\midrule
营业收入（亿元） & 159.4 & 218.6 & 253.8 & 287.0 \\
归母净利润（亿元） & 26.5 & 35.2 & 39.2 & 41.0 \\
毛利率 (\%) & 37.8 & 37.0 & 36.5 & 36.0 \\
净利率 (\%) & 16.6 & 16.1 & 15.4 & 14.3 \\
ROE (\%) & 10.5 & 12.1 & 12.3 & 11.8 \\
\bottomrule
\end{tabularx}
\sourcenote{公司财报，本研究团队预测}
\end{exhibitbox}

公司盈利能力稳健，净利率长期维持在 15\% 左右，在国内功率半导体公司中处于领先水平。轨道交通业务提供稳定现金流，功率半导体业务提供高增长弹性。

\subsection{估值分析}

截至 2026 年 6 月，公司市值约 804 亿元，PE(TTM) 约 19.6 倍，是所有覆盖公司中估值最低的。相对于功率半导体行业平均 60–80 倍的 PE，时代电气估值折价显著，主要原因是市场将其视为轨道交通装备公司而非半导体公司。

我们认为，随着功率半导体业务占比提升，市场将逐步认识到公司的半导体属性，估值有望向半导体行业靠拢，存在显著的估值修复空间。

\subsection{催化剂与风险}

催化剂：
\begin{itemize}[leftmargin=2em]
\item SiC 车规项目大规模量产交付
\item 功率半导体业务分拆上市预期
\item 轨道交通业务订单超预期
\item 半导体业务占比提升带动估值重估
\end{itemize}

风险：
\begin{itemize}[leftmargin=2em]
\item 轨道交通投资不及预期
\item 车规 IGBT/SiC 客户拓展慢于预期
\item 功率半导体业务价格竞争加剧
\end{itemize}

\subsection{投资评级：买入}

我们给予时代电气\textbf{买入}评级。公司是国内功率半导体 IDM 的龙头之一，技术实力突出，车规级进展领先，但估值仍按传统装备公司定价，存在显著的估值修复空间。预计 2025–2027 年净利润 CAGR 约 15–20\%，当前估值具备较高安全边际。

\section{斯达半导（603290）— IGBT 国产替代龙头}

\subsection{基本信息与业务概览}

斯达半导是国内 IGBT 领域的龙头企业，专注于 IGBT 芯片设计和模块封装，是国内第一家实现 IGBT 大规模产业化的公司。公司采用 Fabless 为主、自建部分封装产能的模式，产品覆盖 600V–6500V 全电压等级。

\subsection{业务结构与技术进展}

斯达半导以 IGBT 模块为核心产品，下游应用以工业控制和新能源为主，近年来新能源汽车业务快速增长，已成为第一大下游。

\begin{exhibitbox}[表 6-3：斯达半导下游结构（估算）]
\centering
\small
\begin{tabularx}{\textwidth}{X c c c}
\toprule
\textbf{下游应用} & \textbf{营收占比} & \textbf{同比增速} & \textbf{代表客户} \\
\midrule
新能源汽车 & ~30-35\% & 高速增长 & 多家头部车企 \\
工业控制 & ~35-40\% & 稳步增长 & 汇川、英威腾等 \\
新能源发电（光伏/储能） & ~15-20\% & 快速增长 & 头部逆变器厂商 \\
变频家电 & ~10-15\% & 平稳 & 美的、格力等 \\
其他 & ~5\% & — & — \\
\bottomrule
\end{tabularx}
\sourcenote{公司财报，本研究团队测算}
\end{exhibitbox}

技术进展：公司第七代车规级 IGBT 已批量装车，第八代 IGBT 芯片正在研发中；SiC MOSFET 已推出样品并在客户端验证。

\subsection{AI/数据中心业务弹性}

斯达半导在数据中心领域的布局相对较少，主要产品为工业级 IGBT，应用于服务器 PSU 和 UPS 等场景。AI 服务器电源需求增长对公司有一定正面影响，但弹性相对有限。

\subsection{财务数据与盈利能力}

\begin{exhibitbox}[表 6-4：斯达半导核心财务数据]
\centering
\small
\begin{tabularx}{\textwidth}{X c c c c}
\toprule
\textbf{指标} & \textbf{2022} & \textbf{2023} & \textbf{2024} & \textbf{2025E} \\
\midrule
营业收入（亿元） & 27.1 & 30.5 & 33.9 & 42.0 \\
归母净利润（亿元） & 7.3 & 6.0 & 5.1 & 6.5 \\
毛利率 (\%) & 41.5 & 38.2 & 36.5 & 37.5 \\
净利率 (\%) & 27.0 & 19.6 & 15.0 & 15.5 \\
ROE (\%) & 15.5 & 12.2 & 9.6 & 10.8 \\
\bottomrule
\end{tabularx}
\sourcenote{公司财报，本研究团队预测}
\end{exhibitbox}

公司毛利率和净利率在国内同行业中处于较高水平，反映了其较好的产品结构和技术壁垒。但 2023–2024 年受行业周期和价格竞争影响，盈利能力有所下滑，预计 2025 年起逐步企稳回升。

\subsection{估值分析}

截至 2026 年 6 月，公司市值约 314 亿元，PE(TTM) 约 95 倍，估值处于行业中上水平。相对于其行业地位和成长前景，我们认为估值基本合理但不算便宜。

\subsection{催化剂与风险}

催化剂：SiC 器件大规模量产、车规客户突破超预期、行业周期上行带来的业绩弹性。

风险：IGBT 价格竞争超预期、Fabless 模式下产能供应不稳定、SiC 进展慢于预期。

\subsection{投资评级：增持}

我们给予斯达半导\textbf{增持}评级。公司是国内 IGBT 设计龙头，品牌和客户基础扎实，车规级突破领先。当前估值反映了较高的成长预期，建议逢低配置。

\section{华润微（688396）— 产品线最全功率 IDM 龙头}

\subsection{基本信息与业务概览}

华润微是国内产品线最全的功率半导体 IDM 厂商，拥有 6 英寸、8 英寸、12 英寸多条晶圆产线，产品覆盖 MOSFET、IGBT、二极管、功率 IC 等全系列功率半导体产品。

公司采用"IDM + 代工"双轮驱动模式，既有自有品牌产品，也为 Fabless 公司提供晶圆制造服务。

\subsection{业务结构与技术进展}

\begin{exhibitbox}[表 6-5：华润微业务结构]
\centering
\small
\begin{tabularx}{\textwidth}{X c c c}
\toprule
\textbf{业务板块} & \textbf{营收占比} & \textbf{毛利率} & \textbf{备注} \\
\midrule
功率器件 & ~60\% & ~28-30\% & MOS/IGBT/二极管全系列 \\
集成电路（功率IC） & ~10-15\% & ~35-40\% & 电源管理、驱动IC等 \\
晶圆代工 & ~20-25\% & ~20-25\% & 功率代工为主 \\
其他（设计服务/封测） & ~5-10\% & — & — \\
\midrule
合计 & 100\% & ~28-30\% & — \\
\bottomrule
\end{tabularx}
\sourcenote{公司财报，本研究团队测算}
\end{exhibitbox}

技术进展：公司在低压 MOS 领域国内领先，中高压 MOS 和超结 MOS 持续推进；IGBT 单管和模块逐步突破，车规级 IGBT 进入验证阶段；SiC 二极管已量产，SiC MOSFET 研发推进中。

\subsection{AI/数据中心业务弹性}

华润微在数据中心电源领域有较广泛的布局，产品包括高压 MOS、超结 MOS、PFC 二极管等，受益于 AI 服务器功率提升趋势。公司的 12 英寸产线也为数据中心客户提供代工服务，AI 弹性较好。

\subsection{财务数据与盈利能力}

\begin{exhibitbox}[表 6-6：华润微核心财务数据]
\centering
\small
\begin{tabularx}{\textwidth}{X c c c c}
\toprule
\textbf{指标} & \textbf{2022} & \textbf{2023} & \textbf{2024} & \textbf{2025E} \\
\midrule
营业收入（亿元） & 100.6 & 95.9 & 101.2 & 120.0 \\
归母净利润（亿元） & 26.2 & 10.0 & 7.6 & 10.5 \\
毛利率 (\%) & 38.8 & 28.0 & 24.9 & 27.5 \\
净利率 (\%) & 26.0 & 10.4 & 7.5 & 8.8 \\
ROE (\%) & 15.8 & 5.7 & 4.0 & 5.3 \\
\bottomrule
\end{tabularx}
\sourcenote{公司财报，本研究团队预测}
\end{exhibitbox}

公司 2023–2024 年盈利能力大幅下滑，主要受行业周期下行和价格竞争影响。随着行业复苏和高端产品占比提升，预计 2025 年起盈利能力将逐步修复。

\subsection{估值分析}

截至 2026 年 6 月，公司市值约 1,065 亿元，PE(TTM) 约 117 倍。由于当前处于盈利周期底部，PE 失真较为严重。参考 PS 估值，公司 PS 约 10 倍，高于国际龙头英飞凌（~3.5 倍），反映了市场对其成长和国产替代空间的预期。

\subsection{催化剂与风险}

催化剂：12 英寸产线爬坡、IGBT 车规突破、SiC 进展、行业复苏。

风险：产能扩张消化不及预期、价格竞争加剧、高端产品突破慢于预期。

\subsection{投资评级：增持}

我们给予华润微\textbf{增持}评级。公司是国内产品线最全的功率 IDM，制造平台优势显著，受益于国产替代和行业复苏双重逻辑。当前估值较高但考虑到业绩弹性较大，建议积极关注。

\section{士兰微（600460）— IDM 模式 12 英寸进军}

士兰微是国内 IDM 模式的代表性企业，产品覆盖功率分立器件、功率 IC、MEMS 传感器等多个领域。公司近年大力推进 12 英寸产线建设，战略重心向中高端功率器件倾斜。

业务亮点：12 英寸产线持续爬坡，IGBT 模块在工控和白电领域快速突破，SiC 二极管实现量产。财务方面，2024 年受行业周期和新产线折旧影响，盈利承压。

我们给予士兰微\textbf{中性}评级。公司长期战略清晰，但短期盈利压力较大，估值偏高（PE 约 164 倍）。建议等待盈利能力拐点确认后再行配置。

\section{扬杰科技（300373）— 分立器件龙头 + SiC 后起之秀}

扬杰科技是国内功率分立器件龙头，产品以二极管、整流桥、MOSFET 为主，近年来大力布局 SiC 器件，成为 SiC 领域的后起之秀。

核心亮点：
\begin{itemize}[leftmargin=2em]
\item 二极管和整流桥业务国内领先，现金流稳定
\item SiC 二极管进展迅速，已成为国内重要供应商
\item SiC MOSFET 研发推进中，预计 2026–2027 年量产
\item 盈利能力优秀，净利率长期维持在 15\% 以上
\end{itemize}

估值方面，公司 PE(TTM) 约 55 倍，在行业中处于中等水平。考虑到 SiC 业务的高成长性和确定性，估值相对合理。

我们给予扬杰科技\textbf{增持}评级。公司分立器件基本盘稳固，SiC 业务弹性大，是攻守兼备的优质标的。

\section{三安光电（600703）— 化合物半导体全产业链}

三安光电是国内化合物半导体龙头，业务涵盖 LED 芯片、射频 GaAs、功率 GaN、SiC 等多个领域，是国内布局最全面的化合物半导体公司。

公司面临的挑战：LED 业务竞争激烈，盈利能力承压；化合物半导体业务处于投入期，整体尚未盈利。

我们给予三安光电\textbf{中性}评级。公司技术布局全面，长期战略价值显著，但短期盈利能力和估值存在压力。适合风险偏好较高的投资者长期布局。

\section{天岳先进（688234）— SiC 衬底龙头}

天岳先进是国内 SiC 衬底龙头企业，专注于 SiC 单晶衬底的研发、生产和销售，是国内少数能量产 6 英寸 SiC 衬底的公司。

公司核心看点：
\begin{itemize}[leftmargin=2em]
\item 国内 SiC 衬底技术领先企业，6 英寸导电型衬底已量产
\item 8 英寸衬底研发推进中，有望保持国内领先
\item 全球 SiC 衬底需求快速增长，公司充分受益
\end{itemize}

风险：SiC 衬底价格快速下降压缩毛利率；产能消化风险；8 英寸技术突破慢于预期；海外厂商降价竞争。

估值方面，公司 PS 约 50 倍，显著高于海外 Wolfspeed（~3 倍 PS），估值溢价主要来自高成长预期和国产替代逻辑。

我们给予天岳先进\textbf{中性}评级。公司技术地位突出，行业空间大，但估值较高且尚未盈利，适合高风险偏好投资者。

\section{新洁能（605111）— MOSFET 设计龙头}

新洁能是国内 MOSFET 设计龙头，采用 Fabless 模式，产品覆盖中低压 MOS、超结 MOS、屏蔽栅 MOS 等多个系列。

公司亮点：
\begin{itemize}[leftmargin=2em]
\item 国内 MOSFET 品类最齐全的设计公司之一
\item 下游应用广泛，消费电子、工业、新能源均衡布局
\item 与华虹、中芯国际等代工厂合作紧密
\end{itemize}

AI 弹性方面，公司产品应用于电源适配器、TV 电源、工业电源等领域，数据中心和服务器电源业务有一定占比，受益于 AI 电源升级。

估值 PE(TTM) 约 83 倍，处于行业中等水平。

我们给予新洁能\textbf{增持}评级。公司在 MOSFET 设计领域竞争力突出，品类持续拓展，受益于国产替代和 AI 数据中心双重逻辑。

\section{东微半导（688261）— 高压 MOS 技术领先}

东微半导是国内高压超结 MOS 领域的技术领先企业，专注于高压功率器件的研发和销售。公司在超结 MOS 领域具有较强的技术实力，部分产品性能达到国际先进水平。

核心亮点：
\begin{itemize}[leftmargin=2em]
\item 超结 MOS 技术国内领先，产品性价比高
\item 受益于充电桩、光伏储能、数据中心电源等需求增长
\item 估值相对合理，PE 约 30–40 倍
\end{itemize}

公司在高压超结 MOS 领域的技术积累深厚，AI 服务器 PSU 和 PFC 需求增长对公司有积极影响。

我们给予东微半导\textbf{增持}评级。公司技术实力突出，估值在行业中相对较低，安全边际较好。

\section{晶丰明源（688368）— 电源管理芯片龙头}

晶丰明源是国内电源管理芯片（PMIC）龙头，专注于 AC-DC 电源管理芯片的研发和销售，产品广泛应用于照明、家电、快充、数据中心等领域。

公司亮点：
\begin{itemize}[leftmargin=2em]
\item 国内 AC-DC PMIC 龙头，市场份额领先
\item 受益于 AI 服务器电源需求增长
\item 产品从照明向数据中心、工业电源等高端领域拓展
\end{itemize}

估值 PE(TTM) 约 247 倍，估值较高，主要反映了 AI 数据中心业务的高增长预期。

我们给予晶丰明源\textbf{中性}评级。公司在 PMIC 领域地位突出，AI 弹性大，但当前估值已经反映了较高预期，需等待业绩兑现。

\begin{keyinsight}[重点公司投资地图]
综合 10 家公司分析，我们建议的配置思路是：
\begin{enumerate}[leftmargin=2em]
\item \textbf{核心底仓}：时代电气（估值洼地 + IDM + 车规）、扬杰科技（盈利稳健 + SiC 弹性）
\item \textbf{成长进攻}：斯达半导（IGBT 龙头）、华润微（全品类 IDM）、新洁能（MOS 设计龙头）
\item \textbf{主题配置}：天岳先进（SiC 衬底）、三安光电（化合物半导体全产业链）
\item \textbf{价值弹性}：东微半导（技术领先 + 估值合理）
\end{enumerate}
\end{keyinsight}
